\section{The six-point Cox ring and the precise representation gate}
\label{sec:cox-six-dictionary}

Put $W=\mathfrak h_m$, let
$\ell_i=e_i-m^{-1}\mathbf1$, and let
$X_m=\operatorname{Bl}_{\{[\ell_1],\ldots,[\ell_m]\}}\mathbb P(W)$,
where $m\ge4$.  Write $H$ for the hyperplane class and $E_i$ for
the exceptional divisors.  The points form a projective frame:
their sole linear relation, up to scale, is $\sum_i\ell_i=0$.
Here $\mathbb P(W)$ has dimension $m-2$ and there are $m$ points.

\begin{proposition}[PNEW-011: Cox grading and exact $S_m$ module]
Choose the exceptional sections $z_i$ with their literal
permutation action.  The Picard-graded, $S_m$-equivariant Cox
ring is the Pl\"ucker ring of $\operatorname{Gr}(2,m+1)$, with
\[
 p_{i,m+1}\longmapsto z_i,
 \qquad
 p_{ij}\longmapsto
 \frac{q_i-q_j}{\prod_{h\notin\{i,j\}}z_h},
 \qquad \sum_iq_i=0,
 \tag{PNEW-C1}
\]
and degrees
\[
 \deg p_{i,m+1}=E_i,
 \qquad \deg p_{ij}=H-\sum_{h\notin\{i,j\}}E_h.
 \tag{PNEW-C2}
\]
The group $S_m$ permutes the first $m$ indices and fixes the
last; the convention $p_{ji}=-p_{ij}$ is part of this action.

For integers $c,k\ge0$, set
\[
 a=c-k,\qquad d=(m-2)c-mk,\qquad q=(m-1)c-mk=c+d.
 \tag{PNEW-C3}
\]
If $a<0$ or $d<0$, then
$H^0(X_m,\mathcal O(cH-k\sum E_i))=0$.
Otherwise there is an actual $S_m$-equivariant isomorphism
\[
 H^0\!\left(X_m,\mathcal O(cH-k\sum E_i)\right)
 \simeq
 \left(S_{(q,c)}\mathbb Q^m\right)_{(a^m)}.
 \tag{PNEW-C4}
\]
In particular, for every $\nu\vdash m$, its $[\nu]$ multiplicity
is exactly $p_{\nu,(a)}^{(q,c)}$.
\end{proposition}
\begin{proof}
The ring identification is the $s=1$, $n=m-2$ case of
Castravet and Tevelev (2006, Remark 3.9).  To fix the group
action and all degrees, we give the direct comparison that
extends the fully read inherited A-028 proof.

Take $R=\mathbb Q[x_1,\ldots,x_m,y_1,\ldots,y_m]$ with
$\mathbb G_a$ acting by $y_i\mapsto y_i+t x_i$.  Its invariants
are generated by $x_i$ and $x_i y_j-x_j y_i$, and form the
Pl\"ucker ring with $p_{i,m+1}=x_i$.  One way to verify the full
assertion is to adjoin a last column to a $2\times m$ matrix.
The stabilizer in $SL_2$ of the last column $(0,1)^T$ is this
additive group.  Additively invariant functions extend
uniquely to $SL_2$-invariant functions on the open set where
the last column is nonzero, using either coordinate chart
to choose an $SL_2$ matrix.  The choices differ by the
additive stabilizer.  The complement has codimension two in
affine matrix space, so the regular function extends uniquely
over it.  Restriction is the inverse.  The Cauchy decomposition
of the polynomial algebra shows that the $SL_2$ invariants
of degree $2h$ are $S_{(h,h)}\mathbb Q^{m+1}$.  Products of
minors form a $GL_{m+1}$-stable space containing a nonzero
highest vector of each such summand, hence generate all of it.
Their equations are the Pl\"ucker equations, for example by
the usual two-row straightening basis.  This proves the
stated invariant ring.  Every map commutes with permutations
of the first $m$ columns.

On $u=\prod_i x_i\ne0$ put
\[
 r_i=y_i/x_i,
 \qquad q_i=u\left(\frac1m\sum_jr_j-r_i\right).
 \tag{PNEW-C5}
\]
These $q_i$ are polynomial invariants of sum zero.  The
localized invariant ring is
\[
\mathbb Q[x_i^{\pm1},q_1,\ldots,q_m]/(\sum_iq_i),
\]
since the differences of the $r_i$ are affine coordinates
for the additive quotient.  Moreover
\[
 x_i y_j-x_j y_i
       =\frac{q_i-q_j}{\prod_{h\notin\{i,j\}}x_h}.
 \tag{PNEW-C6}
\]
Assign $\deg x_i=E_i$ and
$\deg y_i=H-\sum_{h\ne i}E_h$, so each $q_i$ has degree $H$.
A homogeneous element after localization is therefore of
the form $f(q)\prod_i x_i^{v_i}$, with $f$ homogeneous.

We check polynomial extension without assuming an abstract
equivariant Cox isomorphism.  Along $x_i=0$, on the locus
where $y_i$ and the other $x_j$ are nonzero, $q$ is a
nonzero scalar multiple of $\ell_i$.  The local projective
coordinates transverse to $[\ell_i]$ are ratios of differences
of the $q_j$.  Their leading terms after substitution are
$x_i(r_j-r_h)/y_i$, for $j,h\ne i$.  There are $m-2$
algebraically independent such differences, and the scalar
trivializing a homogeneous $f$ is a unit.  The first nonzero
homogeneous term in a local expansion of $f$ therefore
remains nonzero after substitution.  Thus
\[
 \operatorname{ord}_{x_i}f(q(x,y))
                =\operatorname{mult}_{[\ell_i]}f.
 \tag{PNEW-C7}
\]
Since the $x_i$ are distinct primes, the localized homogeneous
element is polynomial precisely when
$v_i+\operatorname{mult}_{[\ell_i]}f\ge0$ for every $i$.
These are precisely the blow-up chart conditions for sections
of $\mathcal O(cH+\sum_i v_iE_i)$.  Positive exceptional
coefficients introduce their canonical sections, with no
extra poles downstairs: the points have codimension at
least two and cannot support a pole divisor.  This
identifies every Picard component, its multiplication, and
its actual $S_m$ action, proving (PNEW-C1)--(PNEW-C2).

Let $b_{ij}$ count the factors $p_{ij}$, for $i,j\le m$,
and let $t_i$ count the factors $p_{i,m+1}$ in a monomial
of degree $cH-k\sum_iE_i$.  Equating the coefficients
of $H$ and $E_i$ gives
\[
 \sum_{i<j}b_{ij}=c,
 \qquad t_i+\sum_{j\ne i}b_{ij}=c-k=a.
 \tag{PNEW-C8}
\]
Its first $m$ column weights are therefore all $a$; its
last column weight is
$\sum_i t_i=ma-2c=(m-2)c-mk=d$; and its number of
Pl\"ucker factors is $c+d=q$.
Conversely these column weights and total degree recover
the displayed Picard degree.  The degree-$q$ Pl\"ucker
algebra is $S_{(q,q)}\mathbb Q^{m+1}$.  Negative column
weights yield zero, proving the stated zero cases.

For nonnegative $a,d$, restrict this Schur module to
$GL_m\times GL_1$, with the last coordinate as the second
factor.  The one-column branching rule is multiplicity
free: an interlacing weight $\beta$ satisfies
\[
 q\ge\beta_1\ge q\ge\beta_2\ge0,
 \qquad\beta_i=0\ (i\ge3),
\]
and the last weight is $2q-|\beta|$.  Last weight $d$
therefore forces $\beta=(q,q-d)=(q,c)$, uniquely.
This is a $GL_m$-equivariant identification, so it commutes
with the literal permutation subgroup.  Taking weight
$(a^m)$ proves (PNEW-C4).  Finally the balanced-space
Schur--Weyl identification gives the stated plethysm
multiplicity, with no projective-character ambiguity.
\end{proof}

\subsection{The precise outer-six specialization}
For six points of $\mathbb P^4$ the Grassmannian is
$\operatorname{Gr}(2,7)$.  The dictionary is
\[
 cH-k\sum_{i=1}^6E_i
 \quad\longleftrightarrow\quad
 ((c-k)^6,4c-6k),
 \qquad q=5c-6k.
 \tag{PNEW-C9}
\]
The scalar fat-point module is consequently
$\left(S_{(5c-6k,c)}\mathbb Q^6\right)_{((c-k)^6)}$,
when the nonnegative-weight conditions hold.  In particular
the quadrics through the frame, $c=2,k=1$, form the actual
Specht module $[(4,2)]$, of dimension nine: the weight
$(1^6)$ space of $S_{(4,2)}\mathbb Q^6$ is the usual
Schur--Weyl Specht realization.  Along the ray
$l(2H-\sum E_i)$ the seven column weights are
$(l,l,l,l,l,l,2l)$, so the natural column permutation
symmetry is the first $S_6$; the last column has a different
weight for $l>0$.

\subsection{Why the Cox identification does not finish item 5}
The ordinary Pl\"ucker standard monomial basis counts the
dimension of a component.  It is not a basis permuted by
$S_m$.  For example, the Pl\"ucker relation on four indices
is
\[
 p_{12}p_{34}-p_{13}p_{24}+p_{14}p_{23}=0.
\]
The permutation $(2\ 4)$ sends the standard monomial
$p_{12}p_{34}$ to
$-p_{14}p_{23}=p_{12}p_{34}-p_{13}p_{24}$.
The two terms on the right are standard monomials, with
opposite signs.  More intrinsically, the component
$H^0(X_m,\mathcal O(H))\simeq W$ has trace $-1$ on an
$m$-cycle, precluding any literal permutation basis.
An ordinary toric degeneration or Hilbert basis therefore
does not alone evaluate the individual Specht multiplicities.
Equation (PNEW-C4) makes the remaining finite-group problem
explicit rather than assuming it has disappeared.

There is a second, separate scope issue.  For a general
$\lambda\vdash mn$, G-015 identifies the balanced module
with
\[
 H^0\!\left(\operatorname{Flag}(W),
   \mathcal L_{(\lambda_2,\ldots,\lambda_m)}
       \otimes\pi^*\mathcal I_{Z_m}^{\max(0,\lambda_2-n)}\right).
 \tag{PNEW-C10}
\]
The full tail determines a flag line bundle; pushing it
to $\mathbb P(W)$ generally produces a vector bundle,
not merely $\mathcal O(c)$.  The scalar Cox component above
controls only the corresponding scalar interpolation
problem.  The desired all-$\lambda$, all-$\mu\vdash6$
answer additionally needs the indicated flag/quotient-bundle
sections and their actual $S_6$ covariants.  Neither their
positive multiplicity rule nor an evaluator with the
requested complete $n\le3$ certificate is claimed here.

\paragraph{Source locators and attribution.}
The general Cox identification is known, and is not claimed
as a new theorem of this continuation.  Its primary locator is
Castravet, A.-M., \& Tevelev, J. (2006).
Hilbert's 14th problem and Cox rings.
\emph{Compositio Mathematica, 142}(6), 1479--1498,
Remark 3.9, printed p.~1487.
\url{https://arxiv.org/pdf/math/0505337}.
That remark was independently opened and read; its parameters
are $s=1$, $n=m-2$.  The explicit action and local comparison
above extend A-028, ``The five-point Cox ring and the Segre
cubic, with the exact group action'', in the inherited
\path{sources/RETURN_GCT_PLETH_checkpoint_09_20261121T0546Z/work/paper.md}.
The whole A-028 proof and the whole needed G-015 proof were
read.  The new contribution here is the checked arbitrary-$m$
degree dictionary, the exact balanced-module reduction, and
the stated boundary of what those facts establish for item 5.
